% Revize 2026-10-28: obsah, matematika a návaznost; podrobnosti a časy
% viz planovani/REVIZE_PRIPRAV_1_ROCNIK.md.
% Aktuální porovnání se skutečným obsahem místních PDF Realisticky.cz:
% externi_zdroje/realisticky/01_Zakladni_poznatky/02_Ciselne_obory/010201_Ciselne_obory_prirozena_cisla/010201_Lekce.pdf
% externi_zdroje/realisticky/01_Zakladni_poznatky/02_Ciselne_obory/010202_Cela_cisla/010202_Lekce.pdf
% Rozpis: 1_rocnik/tematicke_plany/Rozpis_vyuky_MAT_1_rocnik.tex, hodina 5.
% Plán: 1_rocnik/tematicke_plany/SPS.pdf, oddíl 1.2.
% Učebnice: ucebnice/01_Zakladni_poznatky/Ciselne_obory/Číselné obory.tex,
% oddíly Přirozená čísla a Celá čísla; návaznost na přípravy 04 a 12.
% Prostudované místní lekce Realisticky.cz (metodická inspirace):
% 01_Zakladni_poznatky/02_Ciselne_obory/010201_Ciselne_obory_prirozena_cisla/010201_Lekce.pdf
% (přehled oborů, Př. 3–6);
% 01_Zakladni_poznatky/02_Ciselne_obory/010202_Cela_cisla/010202_Lekce.pdf
% (Př. 6 a modely znamének);
% 01_Zakladni_poznatky/01_Opakovani/010110_Priority_operaci_zavorky/010110_Lekce.pdf
% (výklad priorit, Př. 4–7). Všechny cesty jsou vzhledem k externi_zdroje/realisticky/.
% Výklad, konkrétní zadání i řešení níže jsou vlastní tvorba.
% Cílená revize: Peanovy axiomy a vlastnosti operací podle zadání učitele.
% Další podklady: v uvedené kapitole učebnice poznámka o Peanových axiomech,
% tabulky vlastností N a Z; Realisticky 010201, oddíl Zajímavé vlastnosti operací,
% a 010202, Př. 1–5. Rozlišena oboustranná neutralita a změna seskupení.
% Nové příklady V1–V3 a R2 jsou vlastní tvorba. Původní zadání zachována.
\documentclass[11pt,a4paper]{article}
% Spolecny styl pro ucitelske pripravy MAT (Work/MAT)
% Verze 1.1, 2026-10-23
\NeedsTeXFormat{LaTeX2e}
\ProvidesPackage{priprava}[2026/10/23 v1.1 Sablona ucitelskych priprav MAT]

% Zakladni typografie a cestina
\RequirePackage[T1]{fontenc}
\RequirePackage[utf8]{inputenc}
\RequirePackage[czech]{babel}
\RequirePackage{lmodern}
\RequirePackage{microtype}
\RequirePackage{amsmath,amssymb}
\RequirePackage{xcolor}
\RequirePackage{enumitem}
\RequirePackage[most]{tcolorbox}
\RequirePackage{tabularx,booktabs,array}
\RequirePackage{fancyhdr}
\RequirePackage{needspace}
\RequirePackage{geometry}
\RequirePackage[hidelinks]{hyperref}

\geometry{a4paper,top=17mm,bottom=18mm,left=23mm,right=23mm}
\definecolor{notegray}{gray}{0.39}
\setlength{\parindent}{0pt}
\setlength{\parskip}{2.5pt plus 1pt minus 0.5pt}
\setlist[itemize]{topsep=2pt,itemsep=1pt,parsep=0pt,leftmargin=1.45em}
\setlist[enumerate]{topsep=2pt,itemsep=2pt,parsep=0pt,leftmargin=1.65em}
\setcounter{secnumdepth}{0}
\renewcommand{\arraystretch}{1.12}

% Zahlavi: vlevo rocnik | globalni tema, vpravo cislo: nazev hodiny
% Zapati zustava zcela prazdne.
\pagestyle{fancy}
\fancyhf{}
\renewcommand{\headrulewidth}{0.35pt}
\renewcommand{\footrulewidth}{0pt}
\setlength{\headheight}{14pt}
\newcommand{\nastavPripravu}[4]{%
  \fancyhead[L]{\small #1\enspace\textbar\enspace #2}%
  \fancyhead[R]{\small #3: #4}%
}

% Minimalisticke cernobile boxy. Kazda definice zvlast.
\tcbset{commonlesson/.style={enhanced,breakable=false,colback=white,colframe=black,
  colbacktitle=white,coltitle=black,boxrule=0.45pt,
  arc=0.8mm,outer arc=0.8mm,left=3mm,right=3mm,top=1.6mm,bottom=1.6mm,
  toptitle=1mm,bottomtitle=0.7mm,fonttitle=\bfseries,
  attach boxed title to top left={xshift=2.5mm,yshift=-1.15mm},
  boxed title style={colback=white,colframe=white,boxrule=0pt,left=1mm,right=1mm}}}
\newtcolorbox{zadaniBox}[1][]{commonlesson,title={Zadání},#1}
\newtcolorbox{prikladBox}[1][]{commonlesson,title={Příklad},#1}
\newtcolorbox{definiceBox}[1][]{commonlesson,title={Definice},#1}

% Spolecny box pro zadani a jeho reseni, bez mezery mezi castmi.
% Pouziti:
%   \begin{prikladReseniBox}
%     Zadani ulohy.
%     \tcblower
%     \begin{align*} ... \end{align*}
%   \end{prikladReseniBox}
% Spodni cast ma jemny sedy podklad a delici linku; lze v ni sazet
% vice radku matematiky. Obe casti maji jen jeden spolecny ramecek.
\tcbset{solutionlesson/.style={
  commonlesson,
  enhanced jigsaw,
  breakable,
  colbacklower=black!3,
  segmentation style={solid,draw=black!25,line width=0.35pt},
  fontlower=\small\color{notegray},
  before lower={\noindent\textbf{Řešení:}\enspace},
  leftlower=3mm,rightlower=3mm
}}
\newtcolorbox{prikladReseniBox}[1][]{solutionlesson,title={Příklad},#1}
\newtcolorbox{zadaniReseniBox}[1][]{solutionlesson,title={Zadání},#1}

% Pomucky pro autora pripravy: sedive poznamky; \reseni pro starsi pripravy
\newcommand{\poznamka}[1]{\par\begingroup\small\color{notegray}\itshape
  \textnormal{\textbf{Poznámka pro učitele.}}\enspace #1\par\endgroup}
\newcommand{\reseni}[1]{\par\begingroup\small\color{notegray}
  \textbf{Očekávané řešení:}\enspace #1\par\endgroup}
\newcommand{\blok}[1]{\needspace{5\baselineskip}\subsection*{#1}}
\newcommand{\vyrok}[1]{\uv{#1}}

\nastavPripravu{1. ročník}{Číselné obory}{05}{Přirozená a celá čísla}
\begin{document}

\poznamka{Rozšířená příprava je zásobníkem pro výuku. Pro 45 minut
orientačně vyber: obory a následník 7 min, vlastnosti operací 15 min,
znaménka a přednost operací 12 min, část samostatné práce 9 min,
shrnutí 2 min. Peanovým axiomům věnuj jen krátkou orientaci; část
procvičování lze přesunout domů. Původní domácí práce je asi na 15 min,
rezervy vybírej podle potřeby, nezadávej je automaticky všechny.}

\blok{Od počítání předmětů k záporným číslům}
Nech žáky vyjádřit počet strojů v dílně, prázdný sklad a teplotu tři stupně
pod nulou. V minulé hodině jsme volili obor proměnné; nyní upřesníme,
která čísla do jednotlivých oborů patří a které výpočty v nich zvládneme.

\begin{definiceBox}
\textbf{Přirozená čísla} jsou čísla $1,2,3,\ldots$; jejich obor značíme
$\mathbb N$. V naší učebnici nulu mezi přirozená čísla nezahrnujeme.
\end{definiceBox}

\begin{definiceBox}
\textbf{Následník} přirozeného čísla $n$ je číslo $n+1$:
například následníkem $4$ je $5$ a následníkem $5$ je $6$.
\end{definiceBox}

\blok{Peanovy axiomy: volitelná zajímavost}
Jak popsat nekonečně mnoho čísel, aniž bychom je všechna vypsali?
Jednou možností jsou \emph{Peanovy axiomy}, tedy výchozí pravidla
pro počáteční prvek a přiřazení následníka $S$. V naší konvenci začínáme
jedničkou. Zápis $S(n)=n+1$ je známý model; při formálním budování
aritmetiky se nejprve zavádí $S$ a sčítání až pomocí něj.
\begin{enumerate}[label=\arabic*)]
\item \textbf{Počáteční prvek:} $1\in\mathbb N$. Máme odkud začít.
\item \textbf{Existence a jednoznačnost následníka:} každému
$n\in\mathbb N$ odpovídá právě jedno $S(n)\in\mathbb N$.
Řetěz může pokračovat a z jednoho čísla se nerozvětvuje do dvou následníků.
\item \textbf{Různá čísla mají různé následníky:}
jestliže $S(m)=S(n)$, pak $m=n$. Dvě různé větve se tedy nespojí.
To je jiný požadavek než jednoznačnost v předchozím bodě.
\item \textbf{Počáteční prvek není následníkem:}
pro žádné $n\in\mathbb N$ neplatí $S(n)=1$.
Nula není protipříklad, protože do našeho $\mathbb N$ nepatří.
\item \textbf{Princip matematické indukce:} jestliže vlastnost $P$
platí pro $1$ a pro každé $n\in\mathbb N$ platí
$P(n)\Rightarrow P(S(n))$, pak $P$ platí pro všechna přirozená čísla.
\end{enumerate}
Indukci si představíme jako nekonečnou řadu domina: potřebujeme
první padající kostku \emph{i} zaručený přechod od každé kostky k následující.
Pouhé ověření prvních sta případů obecný přechod nenahradí.
Formálně indukce také vylučuje další prvky mimo řetěz začínající jedničkou:
množina obsahující $1$ a s každým svým prvkem i jeho následníka
musí obsahovat celé $\mathbb N$.

Například tvrzení \uv{$1+2+\cdots+n=\frac{n(n+1)}2$} má počáteční
případ $1=\frac{1\cdot2}2$. Pokud platí pro některé $n$, pak
\[
1+2+\cdots+n+(n+1)
=\frac{n(n+1)}2+(n+1)=\frac{(n+1)(n+2)}2.
\]
Tím jsme ukázali obecný přechod, ne pouze další jednotlivý výpočet.

V jiné běžné konvenci je počátečním prvkem $0$ a indukce začíná nulou.
Takový obor zde označujeme $\mathbb N_0=\{0,1,2,\ldots\}$.
Při čtení jiného zdroje proto vždy ověříme jeho konvenci.

\poznamka{Jde pouze o intuitivní představení jedné možnosti formálního
základu. Celý exkurz i ukázka indukce jsou nepovinné; nejsou požadavkem
k testu. Při běžné hodině stačí obraz řetězu a dva nutné kroky indukce;
podrobný rozbor může zůstat pro zájemce.
Nech žáky určit následníka $99$ a vysvětlit, proč přirozená čísla nemají poslední prvek.}

\begin{definiceBox}
\textbf{Celá čísla} jsou přirozená čísla, jejich záporné protějšky a nula.
Jejich obor značíme
\[
\mathbb Z=\{\ldots,-3,-2,-1,0,1,2,3,\ldots\}.
\]
\end{definiceBox}

Zápis $a\in\mathbb Z$ znamená, že $a$ je celé číslo; $a\notin\mathbb N$
znamená, že není přirozené. Každé přirozené číslo je celé, opačně to neplatí.
Čísla větší než nula jsou kladná, menší než nula záporná. Nula není ani
kladná, ani záporná; je současně nezáporná i nekladná.

\begin{definiceBox}
\textbf{Opačné číslo} k číslu $a$ je číslo $-a$; jejich součet je nula:
$a+(-a)=0$. K nule je opačným číslem opět nula.
\end{definiceBox}

\begin{prikladReseniBox}[title={P1: číslo a jeho protějšek},breakable=false]
Urči opačná čísla k $7$, $-9$, $0$ a k hodnotě $4-11$.
Je zápis $-a$ vždy záporné číslo?
\tcblower
Postupně dostáváme $-7$, $9$, $0$ a $7$, protože $4-11=-7$.
Předpona minus označuje opačné číslo, nezaručuje zápornost:
pro $a=-9$ je $-a=9$, pro $a=0$ je $-a=0$.
\end{prikladReseniBox}

\blok{Která pravidla smíme při počítání používat?}
Nejprve zvolíme obor a operaci. Například $2+5=7$ zůstává v $\mathbb N$,
ale $2-5=-3$ už ne. Stejná čísla tak mohou při jiné operaci vyžadovat
rozšíření oboru. Jediný protipříklad vyvrátí obecné pravidlo;
několik správných výpočtů samo o sobě obecnou platnost nedokazuje.

\begin{definiceBox}
\textbf{Uzavřenost oboru vzhledem k operaci} znamená, že pro každé dva
prvky tohoto oboru je výsledek dané operace definován a patří opět do téhož oboru.
\end{definiceBox}

Sčítání a násobení jsou uzavřené v $\mathbb N$ i $\mathbb Z$.
Odčítání je uzavřené v $\mathbb Z$, v $\mathbb N$ nikoli:
už $3-3=0$ není podle naší konvence přirozené číslo.
Vlastnosti nejsou tvrzením, že \emph{některý} výpočet vyjde, ale že
vyjdou \emph{všechny} přípustné výpočty daného typu.

\begin{definiceBox}
\textbf{Komutativita} znamená, že záměna pořadí obou čísel nemění výsledek
operace. Pro sčítání píšeme $a+b=b+a$, pro násobení $a\cdot b=b\cdot a$.
\end{definiceBox}

\begin{definiceBox}
\textbf{Asociativita} znamená, že změna seskupení tří čísel nemění
výsledek operace: $(a+b)+c=a+(b+c)$, případně
$(a\cdot b)\cdot c=a\cdot(b\cdot c)$. Pořadí čísel přitom zůstává stejné.
\end{definiceBox}

Sčítání i násobení jsou v $\mathbb N$ a $\mathbb Z$ komutativní
a asociativní. Odčítání v $\mathbb Z$ nemá ani jednu z těchto vlastností.

\begin{prikladReseniBox}[title={V1: pořadí a závorky jsou dvě různé věci},breakable=false]
Porovnej a pojmenuj použitou vlastnost, nebo uveď, proč rovnost neplatí:
\begin{enumerate}[label=\alph*)]
\item $(-4)+9$ a $9+(-4)$;\quad $(-2)\cdot7$ a $7\cdot(-2)$.
\item $(2+5)+3$ a $2+(5+3)$;\quad $(2\cdot5)\cdot3$ a $2\cdot(5\cdot3)$.
\item $8-3$ a $3-8$;\quad $(8-3)-2$ a $8-(3-2)$.
\end{enumerate}
\tcblower
a) Oba součty jsou $5$, oba součiny $-14$; jde o komutativitu v $\mathbb Z$.
b) Součty jsou $10$, součiny $30$; jde o asociativitu, zde už v $\mathbb N$.
c) Dostáváme $5\ne-5$ a $3\ne7$. Jsou to protipříklady komutativity
a asociativity odčítání v $\mathbb Z$; všechny uvedené rozdíly jsou přitom celé.
\end{prikladReseniBox}

\begin{definiceBox}
\textbf{Neutrální prvek} dané operace v daném oboru je jeho prvek,
který při této operaci s libovolným prvkem oboru nezmění jeho hodnotu
v žádném z obou pořadí.
\end{definiceBox}

V $\mathbb Z$ je pro sčítání neutrální $0$: $a+0=0+a=a$.
Pro násobení je v $\mathbb N$ i $\mathbb Z$ neutrální $1$: $a\cdot1=1\cdot a=a$.
V $\mathbb N$ s naší konvencí sčítání neutrální prvek nemá,
protože nula do tohoto oboru nepatří a přičtení kladného čísla hodnotu zvětší.

\begin{prikladReseniBox}[title={V2: neutralita a opačný prvek závisejí na operaci a oboru},breakable=false]
a) Doplň v $\mathbb Z$: $-7+\square=-7$; $-7\cdot\square=-7$;
$-7+\square=0$. b) Je $0$ neutrální pro odčítání?
c) Má číslo $5$ opačný prvek pro sčítání v $\mathbb N$ i $\mathbb Z$?
\tcblower
a) Postupně $0$, $1$, $7$. První dvě čísla zachovají druhý prvek;
třetí je opačný prvek a součet změní na neutrální nulu.
b) Ne: ačkoli $7-0=7$, při obráceném pořadí $0-7=-7\ne7$.
Žádný oboustranný neutrální prvek odčítání v $\mathbb Z$ není:
z požadavku $a-e=a$ plyne $e=0$, které nevyhoví druhému pořadí.
c) V $\mathbb Z$ ano, je to $-5$. V $\mathbb N$ ne, protože $-5\notin\mathbb N$.
Každé celé číslo má opačné celé číslo; žádné přirozené číslo nemá opačný prvek v $\mathbb N$.
\end{prikladReseniBox}

\poznamka{Opačný prvek zde vždy vztahuj ke sčítání: dává s původním
číslem neutrální nulu. Neutrální prvek původní číslo zachová.
Převrácený prvek vzhledem k násobení přijde v hodině 06.
Názvy vlastností žáci vysvětlují na výpočtech; není nutné memorovat
abstraktní seznam bez porozumění.}

\begin{definiceBox}
\textbf{Inverzní prvek} k prvku $a$ vzhledem k dané operaci
je prvek téhož oboru, který s $a$ v obou pořadích dává neutrální prvek
této operace. Pro sčítání v $\mathbb Z$ je to opačné číslo $-a$.
\end{definiceBox}

\blok{Přehled vlastností operací}
V tabulce jsou $a,b,c$ libovolná čísla uvedeného oboru; zápisy
se vztahují vždy k pojmenované operaci. U neutrálního a inverzního prvku
musí hledaný prvek patřit do zvoleného oboru.
\begin{center}
\small
\begin{tabularx}{\linewidth}{@{}>{\raggedright\arraybackslash}p{24mm}>{\raggedright\arraybackslash}p{37mm}>{\raggedright\arraybackslash}X>{\raggedright\arraybackslash}X@{}}
\toprule
Vlastnost & Matematický zápis & Slovní význam & Příklad/protipříklad\\
\midrule
Uzavřenost & $a+b\in\mathbb N$ pro $a,b\in\mathbb N$
& Výsledek vždy zůstane v oboru.
& $2+5=7\in\mathbb N$; odčítání: $2-5\notin\mathbb N$.\\[2mm]
Komutativita & $a+b=b+a$
& Lze zaměnit pořadí.
& $2+5=5+2$; $2-5\ne5-2$.\\[2mm]
Asociativita & $(a+b)+c$\newline ${}=a+(b+c)$
& Lze změnit seskupení při stejném pořadí.
& $(8-3)-2\ne8-(3-2)$.\\[2mm]
Neutrální prvek & $a+0=0+a=a$\newline $a\cdot1=1\cdot a=a$
& Hodnota partnera se nezmění.
& Pro $+$ v $\mathbb Z$: $0$; v našem $\mathbb N$ neexistuje.\\[2mm]
Inverzní prvek & $a+(-a)$\newline ${}=(-a)+a=0$
& Spolu s původním prvkem dává neutrální prvek.
& V $\mathbb Z$: $5+(-5)=0$; $-5\notin\mathbb N$.\\
\bottomrule
\end{tabularx}
\end{center}
Pro násobení platí také komutativita a asociativita. Inverzní prvek
pro násobení by musel dávat součin $1$; například ke $2$ v $\mathbb Z$
neexistuje. K tomu se vrátíme při rozšíření na zlomky.

\begin{zadaniReseniBox}[title={V3: krátké procvičení vlastností},breakable=false]
a) Zdůvodni, proč z $9-4=5$ neplyne uzavřenost odčítání v $\mathbb N$.
b) Zjednoduš výpočet $(-17)+28+17+72$ a pojmenuj změny pořadí a seskupení.
c) Je nula neutrálním prvkem násobení v $\mathbb Z$?
\tcblower
a) Jeden výpočet nestačí; například $4-9=-5\notin\mathbb N$.
b) Komutativita a asociativita sčítání dovolují
$[(-17)+17]+(28+72)=0+100=100$; používáme také opačná čísla a neutrální nulu.
c) Ne, například $6\cdot0=0\ne6$. Nula součin vynuluje, nezachová druhý činitel.
\end{zadaniReseniBox}

\blok{Co znamenají znaménka při počítání}
Odčítání nahradíme přičtením opačného čísla: $a-b=a+(-b)$.
Proto $-6-(-4)=-6+4=-2$, kdežto $-6-4=-10$.
Při násobení a dělení nenulových čísel dávají stejná znaménka kladný
výsledek, různá znaménka záporný. V součinu s nulou je výsledek nula.

Proč součin dvou záporných čísel vyjde kladný? Ze známého roznásobení
dostaneme $0=(-3)\cdot(4+(-4))=-12+(-3)\cdot(-4)$;
poslední součin tedy musí být $12$.

\poznamka{Pravidlo \uv{minus a minus je plus} vztahuj ke konkrétní operaci.
Součet $(-3)+(-4)$ je záporný. Dělení nulou není definováno:
$8:0$ nemá výsledek, protože $0\cdot c$ nikdy není $8$;
u $0:0$ by vyhovovalo každé $c$, takže nevzniká jednoznačný podíl.}

\blok{Závorky určují, co k sobě patří}
Nejprve vypočítáme vnitřní závorky, potom známé mocniny,
pak násobení a dělení, nakonec sčítání a odčítání.
Operace stejné úrovně vyhodnocujeme zleva doprava.
Druhou mocninu zde jen připomínáme ze ZŠ: $3^2=3\cdot3$;
pravidla pro mocniny přijdou později.

\begin{prikladReseniBox}[title={P2: stejná čísla, různý postup},breakable=false]
Vypočti a vysvětli rozdíly:
\[
18:3\cdot2,\qquad 18:(3\cdot2),\qquad -3^2,\qquad (-3)^2.
\]
\tcblower
$18:3\cdot2=6\cdot2=12$, ale $18:(3\cdot2)=18:6=3$.
Násobení nemá přednost před dělením.
Zápis $-3^2=-(3\cdot3)=-9$, zatímco $(-3)^2=(-3)\cdot(-3)=9$.
Závorky rozhodují, zda se násobí také záporné znaménko.
\end{prikladReseniBox}

\begin{prikladReseniBox}[title={P3: více operací v jednom výrazu},breakable=false]
Vypočti $-24:(-6)+3\cdot[2-(-5)]-(-2)^2$.
\tcblower
\begin{align*}
-24:(-6)+3\cdot[2-(-5)]-(-2)^2
 &=4+3\cdot7-4\\
 &=4+21-4=21.
\end{align*}
Znaménko před poslední mocninou znamená odčítání její hodnoty $4$.
\end{prikladReseniBox}

\blok{Samostatná práce: nejprve postup, potom výsledek}
\begin{zadaniReseniBox}[title={S1: základní sada},breakable=false]
U a) a b) urči také, zda je výsledek přirozený, pouze celý, nebo není celý.
\begin{enumerate}[label=\alph*)]
\item $-8+13$;\quad $6-14$;\quad $4-4$;\quad $5:2$.
\item $7-(-3)$;\quad $-4-(-9)$;\quad $(-6)\cdot(-2)$;\quad $0:(-5)$.
\item $36:(-3)\cdot2$;\quad $20-3\cdot(4-7)$.
\item $-2\cdot[5-3\cdot(-2)]$;\quad $(-2)^2-2^2$.
\end{enumerate}
\tcblower
a) $5\in\mathbb N$, $-8\in\mathbb Z$ a $-8\notin\mathbb N$,
$0\in\mathbb Z$ a $0\notin\mathbb N$; $5:2=2{,}5$ není celé číslo.

b) $10,5,12$ jsou přirozená; $0$ je pouze celé.

c) $(-12)\cdot2=-24$; $20-3\cdot(-3)=29$.

d) $-2\cdot(5+6)=-22$; $4-4=0$.
\end{zadaniReseniBox}

\begin{zadaniReseniBox}[title={S2: najdi první chybný krok},breakable=false]
Oprav: a) $-5-(-8)=-13$; b) $24:4\cdot3=24:12=2$;
c) $-4^2=16$. Vysvětli, jaké pravidlo bylo porušeno.
\tcblower
a) Odčítáme záporné číslo: $-5+8=3$.
b) Postupujeme zleva: $6\cdot3=18$.
c) Mocnina má přednost: $-(4\cdot4)=-16$.
\end{zadaniReseniBox}

\poznamka{Zadání dávej bez řešení. Ptej se na první krok, ne pouze na
výsledek. Na závěr nech každého vysvětlit, proč $3-3$ není při naší konvenci
přirozené číslo. Neúplnost oboru $\mathbb Z$ při dělení je mostem k hodině 06.}

\blok{Rezerva pro svižnější práci}
\begin{zadaniReseniBox}[title={R1: volitelně, přibližně 10–15 minut},breakable=false]
\begin{enumerate}[label=\alph*)]
\item Spočti $-\{18-2\cdot[5-(-3)]\}+30:(-5)$.
\item Vlož závorky do $16-8:2$ tak, aby výsledek byl $4$.
Jaká je hodnota původního výrazu?
\item Petr tvrdí: \uv{Pro každé celé $a$ je $a-a$ přirozené číslo.}
Rozhodni. Je součin dvou celých nenulových čísel vždy přirozený?
\item Vysvětli bez dlouhého násobení:
$27\cdot(-19)+27\cdot19+(-4)\cdot25$.
\end{enumerate}
\tcblower
a) $-(18-16)-6=-8$.
b) $(16-8):2=4$; bez závorek $16-4=12$.
c) První tvrzení neplatí: pro každé $a$ je $a-a=0\notin\mathbb N$.
Ani druhé neplatí, například $2\cdot(-3)=-6$.
d) První dva součiny jsou opačná čísla a vyruší se; zbývá $-100$.
\end{zadaniReseniBox}

\poznamka{Rezerva rozšiřuje již známé operace, nezavádí nová pravidla
pro mocniny. Žákům, kteří chybují ve znaménkách, dej místo R1 znovu S1b)
s vysvětlením pomocí opačného čísla.}

\begin{zadaniReseniBox}[title={R2: volitelné doplnění k vlastnostem},breakable=false]
\begin{enumerate}[label=\alph*)]
\item Zdůvodni, proč v $\mathbb N$ existuje neutrální prvek násobení,
ale ne sčítání. Uveď jejich kandidáty.
\item Porovnej $(12-5)-3$ a $12-(5-3)$. Kterou vlastnost tento příklad vyvrací?
\item Vypočti výhodně $(-4)\cdot25\cdot(-3)$ a popiš použité vlastnosti.
\end{enumerate}
\tcblower
a) Jednička patří do $\mathbb N$ a násobení jí číslo nemění.
Pro sčítání by byla nutná nula; ta do našeho $\mathbb N$ nepatří.
b) $4\ne10$; odčítání není asociativní, ani když zde všechny mezivýsledky
i výsledky zůstaly přirozené.
c) $[(-4)\cdot25]\cdot(-3)=(-100)\cdot(-3)=300$.
Měníme seskupení pomocí asociativity násobení v $\mathbb Z$.
\end{zadaniReseniBox}

\poznamka{R2 je další volitelná zásoba, případně náhrada části domácí
práce D1. Po tomto rozšíření není cílem stihnout všechny sady v jedné hodině.}

\blok{Domácí procvičování}
\begin{zadaniReseniBox}[title={D1: přibližně 15 minut},breakable=false]
U výpočtů zapiš alespoň jeden mezikrok.
\begin{enumerate}[label=\arabic*)]
\item Rozhodni o příslušnosti čísel $-1,0,4$ do $\mathbb N$ a $\mathbb Z$.
\item Urči opačné číslo k $-12$, $0$ a k hodnotě $3-8$.
\item Vypočti $-15+7-(-6)$ a $(-7)\cdot(-3)-30$.
\item Porovnej hodnoty $48:6\cdot2$ a $48:(6\cdot2)$.
\item Vypočti $5-2\cdot[3-(-4)]$ a $-5^2+(-5)^2$.
\item Oprav tvrzení: \uv{Jestliže je celé číslo nezáporné, je přirozené.}
Uveď protipříklad a pravdivé znění s co nejmenší změnou.
\end{enumerate}
\tcblower
1. Všechna tři čísla jsou celá, přirozené je pouze $4$.
2. $12,0,5$.
3. $-8+6=-2$; $21-30=-9$.
4. $8\cdot2=16$; $48:12=4$.
5. $5-2\cdot7=-9$; $-25+25=0$.
6. Protipříklad $0$. Správně: každé \emph{kladné celé} číslo je přirozené.
\end{zadaniReseniBox}

\end{document}
